\begin{abstract}
Wasser zu simulieren ist im Bereich der Computergrafik immer noch eine gro"se Herausforderung.
Auch wenn auf dem Gebiet der Offline-Simulation zunehmend plausible Resultate zu sehen sind, wie z.B. in den Kinofilmen "`Titanic"' (1997), "`Antz"' (1998) oder "`Finding Nemo"' (2003), ist das Simulieren von Fl"ussigkeiten im interaktiven Bereich ein weitgehend unbestelltes Feld.
Dank der erreichten Leistungsst"arke aktueller PCs ist eine solche Simulation mit vern"unftigem Aufwand realisierbar.

In der vorliegenden Arbeit soll der Ansatz von Matthias M"uller, David Charypar \& Markus Gross \cite{ETHZ} fortgef"uhrt werden und durch eine neuartige Datenstruktur der Speicherbedarf bei gleichzeitiger Vergr"o"serung des Simulationsgebiets verringert werden.
W"ahrend das Programm zur vorliegenden Arbeit bei h"oheren Bildraten mit unter 15MB RAM auskommt, konnte das Simulationsgebiet auf das etwa hunderttausendfache gesteigert werden.
\end{abstract}